\documentclass[10pt,a4paper]{amsart}
\usepackage[margin=19mm]{geometry}
\usepackage{amsmath,amssymb,mathtools}
\usepackage[scr=rsfs,cal=euler]{mathalfa}
\usepackage{xfrac}
\usepackage[unicode,hidelinks,bookmarks=false]{hyperref}
\newcommand{\ie}{i.e.\ }
\newcommand{\cf}{cf.\ }
\newcommand{\resp}{respectively\ }
\newcommand{\Subsection}{\subsection}
\providecommand{\nobreakdash}{\nobreak\hbox{-}}
\setlength{\emergencystretch}{2em}
\multlinegap0pt
\usepackage{xcolor}
\usepackage{amssymb}
\usepackage{bm}
\newtheorem{lemma}{Lemma}
\title{Global well-posedness in the critical Besov space of the skew mean curvature flow in $\mathbb{R}^d: d\ge 5$}
\author{Ning-An Lai, Jie Shao, Zexian Zhang, Yi Zhou}
\date{Source: arXiv:2505.03609v3 (CC BY 4.0)}
\begin{document}
\maketitle

\noindent\emph{Source and licence.} Global well-posedness in the critical Besov space of the skew mean curvature flow in $\mathbb{R}^d: d\ge 5$. Version arXiv:2505.03609v3 (CC BY 4.0), distributed by arXiv under the Creative Commons Attribution 4.0 International licence, http://creativecommons.org/licenses/by/4.0/, as recorded in the arXiv licence metadata for this exact version and shown on the abstract page https://arxiv.org/abs/2505.03609v3. Full licence notice: \texttt{licenses/arxiv-2505.03609-CC-BY-4.0.txt}.\par\smallskip

\noindent\textit{Editorial scope.} Counted unit: the article's lemma lem\_J11t4 (source0.tex lines 1579--1586). The article's complete original proof of it is delivered with it (source0.tex lines 1587--1807, one \texttt{proof} environment of 221 lines). No mathematics is written, repaired or re-derived: the statement and the proof are the article's own lines, in the article's own notation, and no retained line is substituted or re-flowed.  Retained uncounted as the source-local context the proof needs: lines 311--325; lines 359--364; lines 721--728; lines 1483--1490.  Nothing was omitted from any retained span, so the package carries no omission record.  No retained line contains a citation, so the wrapper carries no bibliography and no bibliography entry.  No retained line contains a figure environment, an image-inclusion command or drawing code, so no image is delivered and the package declares no asset dependency.  Editorial formatting only, in addition: an AMS \texttt{amsart} preamble that restates the article's own declarations and macros used by the retained lines, namely the environments \texttt{lemma} on separate counters, together with the standard packages those lines need.  The retained environments print in this document as Lemma 1 in order of appearance, whereas the article prints the same items with the numbering of its own full document.  The complete native source of the article as distributed in the arXiv e-print for this exact version is bundled unchanged as \texttt{sources/177681/source0.tex}.
\begin{align}
R_{\sigma\gamma\alpha\beta}
=&\,
\operatorname{Re}
\left(
\lambda_{\beta\gamma}\overline{\lambda}_{\alpha\sigma}
-
\lambda_{\alpha\gamma}\overline{\lambda}_{\beta\sigma}
\right),
\label{eq:gauss-complex}\\
\nabla_\alpha^A\lambda_{\beta\gamma}
=&\,
\nabla_\beta^A\lambda_{\alpha\gamma},
\label{eq:codazzi-complex}
\end{align}

\begin{align}
   (\lambda\wedge\lambda)_{\gamma\sigma}:= \lambda_{\gamma\sigma}\overline\psi
-
\lambda_{\alpha\gamma}
\overline{\lambda}^\alpha_{\sigma}, \label{wed.lam}
\end{align}

\begin{align}
	&\|P_{\rho} u_i(t) \|_{L^2_x} \lesssim  \varepsilon^{-\frac18}\rho^{-s-\sigma} c_\rho (\sigma),\label{sup1}\\
	&\left\|P_{\rho }  u_{\mathfrak{T}}\right\|_{L^2_{t,x } }\lesssim   \varepsilon^{-\frac 18} \rho^{-s-1-\sigma} c_\rho(\sigma),\label{supV} \\
	\sup_{\tau}&\,  \big\| \|P_\mu u_i\|_{L^2_{\widehat{x}_\gamma}}   \  \| P_{\rho,\gamma} u_j^\tau\|_{L^2_{\widehat{x}_\gamma}}  \big\|_{L^2_{t,x_\gamma}} \label{sup2} \\
 &\quad \lesssim \varepsilon^{-\frac14} \mu^{-s} c_\mu(0)   \rho^{-s-\frac 12-\sigma} c_\rho(\sigma), ~	\text{ for } \mu\ll \rho,\nonumber\\
			\sup_J &\, \|  P_\rho u_i \|_{L^4_{t,y_J} L^2_{z_J}}  \lesssim \varepsilon^{-\frac18} \rho^{-s+\frac 14-\sigma} c_\rho(\sigma), \label{sup3}\\
	&\left\|  B_{\rho }\right\|_{L^2_{t,x } }\lesssim  \varepsilon^{-\frac 18} \rho^{-s-\sigma} c_\rho(\sigma)\label{supB2},
\end{align}

    \begin{align}
        &  \sum_{\mu\ll \rho}  \sum_{\gamma=1}^d \mu^{-\frac 12} \big\| \left\|P_{\mu } u_i\right\|_{L^\infty_{\widehat{x}_\gamma}} \left\|P_{  \rho,\gamma}    u_j\right\|_{L^2_{\widehat{x}_\gamma}}\big\|_{L^2_{t,x_\gamma}} \lesssim  
 \varepsilon^{\frac 34}     \rho^{-s-\frac12-\sigma} c_{\rho} (\sigma), \label{pre.bi1}\\
 & \sum_{\mu \ll \rho} \mu^{-1} \rho \sup_\tau\Big\| \|    	L(P_{ \mu} u_i, P_{\ll \mu }u_j)  \|_{L^\infty_z} \|    \psi_{\sim \rho,1} \|_{L^2_z}   \|     \psi _{\rho,1}^\tau\|_{L^2_z}     \Big\|_{L^1_{t,y}}   \nonumber  \\
 \lesssim&\, \varepsilon^{\frac32}    \rho^{-2s  -2\sigma}c_{\rho }^2 (\sigma), \label{pre.bi2} \\
 & \sum_{\mu \ll \rho} \mu^{-\frac 32} \sup_\tau  \Big\| \|    	 P_{  \mu } (u_i u_j) \|_{L^\infty_z} \|   P_{ \rho,1} u_l^\tau \|_{L^2_z}     \Big\|_{L^2_{t,y}}   \nonumber  \\
 \lesssim&\, \varepsilon^{\frac 32} \rho^{-s-\frac12-\sigma} c_\rho(\sigma), \label{pre.bi3} 
    \end{align}

 \begin{lemma}  \label{lem_J11t4}
Let \(d\geq5\) and \(y_1=x_1\) in \(y=(y_1,y_2,y_3)\). Then
\begin{align} 
 		& \sum_{\mu \ll \rho}  \rho \mu^{-1} \sup_\tau \big\| \| P_{\mu} \left(\lambda\wedge\lambda\right) \|_{L^\infty_z} \|     \psi_{\sim \rho ,1}\|_{L^2_z}   \|     \psi_{\rho,1}^\tau\|_{L^2_z}     \big\|_{L^1_{t,y}}  \nonumber\\
 &\lesssim  \varepsilon^{\frac32} \rho^{-2s -2\sigma}  c_{\rho }^2 (\sigma),\label{est_J11t4}
 	\end{align}
    where \(\lambda\wedge\lambda\) is defined in \eqref{wed.lam}.
 \end{lemma}

 \begin{proof}
Applying Bony's decomposition to \(\lambda\wedge\lambda\), we obtain
\begin{align}
P_\mu(\lambda\wedge\lambda)
=&\,
L\bigl(P_{\ll\mu}\lambda,P_\mu\lambda\bigr)
+
L\bigl(P_\mu\lambda,P_{\ll\mu}\lambda\bigr)  
\nonumber\\
&+ \sum_{\mu\lesssim \nu\sim\nu'\ll\rho} P_\mu L(P_\nu\lambda,P_{\nu'}\lambda)+
\sum_{\nu\sim\nu'\gtrsim\rho}
P_\mu\bigl(\lambda_\nu\wedge\lambda_{\nu'}\bigr).
\label{gauss.bony}
\end{align}


Plugging the first two terms into \eqref{est_J11t4}, we can see it has been bounded by \eqref{pre.bi2}. The third term can be estimated by using \eqref{sup2} in a way similar  to those for \eqref{pre.bi2} as follows:
\begin{align*}
&\sum_{\mu\ll\rho}
\sum_{\mu\lesssim\nu\sim\nu'\ll\rho}
\rho\mu^{-1}
\sup_\tau
\left\|
\|P_\mu L(P_\nu\lambda,P_{\nu'}\lambda)\|_{L_z^\infty}
\|\psi_{\sim\rho,1}\|_{L_z^2}
\|\psi_{\rho,1}^\tau\|_{L_z^2}
\right\|_{L_{t,y}^1}
\\
\lesssim&\, \sum_{\mu\ll\rho}
\sum_{\mu\lesssim\nu\sim\nu'\ll\rho}
\rho\mu^{d-2}
\sup_\tau
\left\|
\|  L(P_\nu\lambda,P_{\nu'}\lambda)\|_{L_{\widehat{x}_1}^1}
\|\psi_{\sim\rho,1}\|_{L_{\widehat{x}_1}^2}
\|\psi_{\rho,1}^\tau\|_{L_{\widehat{x}_1}^2}
\right\|_{L_{t,x_1}^1}\\
\lesssim &\,
\varepsilon^{\frac32}
\rho^{-2s-2\sigma}
c_\rho^2(\sigma).
\end{align*}

As for the fourth term in \eqref{gauss.bony}, we  need to exploit the additional structure  given by \eqref{wed.lam}. For each dyadic frequency \(\nu\), write
\begin{align}
P_\nu
=
\nu^{-1}
\sum_{\beta=1}^d
\partial_\beta\widetilde P_{\nu,\beta},
\label{Pnu.derivative}
\end{align}
where the multipliers \(\widetilde P_{\nu,\beta}\) are uniformly bounded.
By \(u_3=\partial_xh\), the Codazzi equations \eqref{eq:codazzi-complex} imply schematically
\begin{align}
\partial_\beta\lambda_{\gamma\sigma}
&=
\partial_\sigma\lambda_{\gamma\beta}
+
L\bigl((u_3,A),\lambda\bigr),
\label{codazzi.one}
\\
\partial_\alpha\lambda^\alpha_\sigma
&=
\partial_\sigma\psi
+
L\bigl((u_3,A),\lambda\bigr),
\label{codazzi.contracted}
\end{align}
where we write \( (u_3,A)=u_3+A\) for convenience.

For \(\nu\sim\nu'\gtrsim\rho\), we obtain
\begin{align*}
P_\mu(\lambda_\nu\wedge\lambda_{\nu'})
=&\,
\nu^{-1}P_\mu
\left[
\partial_\sigma
\bigl(\widetilde P_{\nu,\beta}\lambda_{\gamma\beta}\bigr)
\overline{\psi_{\nu'}}
-
\partial_\alpha
\bigl(\widetilde P_{\nu,\beta}\lambda_{\beta\gamma}\bigr)
(P_{\nu^\prime}\overline{\lambda^\alpha _{\sigma}})
\right]
\\
&+
\nu^{-1}P_\mu
L\left(
P_\nu[\lambda\cdot(u_3,A)],
P_{\nu'}\lambda
\right).
\end{align*}
Applying the product rule and using
\(\lambda_{\gamma\beta}=\lambda_{\beta\gamma}\), we further get
\begin{align*}
P_\mu(\lambda_\nu\wedge\lambda_{\nu'})
=&\,
\frac{\mu}{\nu}
L(P_\nu\lambda,P_{\nu'}\lambda)
+
\nu^{-1}P_\mu
\left[
\widetilde P_{\nu,\beta}\lambda_{\beta\gamma}
\left(
\partial_\alpha P_{\nu^\prime}\overline{\lambda^\alpha _{\sigma}}
-
\partial_\sigma\overline{\psi_{\nu'}}
\right)
\right]
\\
&+
\nu^{-1}P_\mu
L\left(
P_\nu[\lambda\cdot(u_3,A)],
P_{\nu'}\lambda
\right),
\end{align*}
which combining the contracted Codazzi equation \eqref{codazzi.contracted}
yields
\begin{align}
P_\mu(\lambda_\nu\wedge\lambda_{\nu'})
=&\,
\frac{\mu}{\nu}
L(P_\nu\lambda,P_{\nu'}\lambda)+
\nu^{-1}P_\mu
L\left(
P_\nu[\lambda\cdot(u_3,A)],
\lambda_{\nu'}
\right)
\label{gauss.high.high.identity}
\\
&
 +
\nu^{-1}P_\mu
L\left(
P_{\nu'}[\lambda\cdot(u_3,A)],
\lambda_\nu
\right).
\nonumber
\end{align}
Thus the quadratic high--high interaction gains the factor
\(\mu/\nu\), while the remaining terms are cubic.

Putting  the first term in right hand side of in
\eqref{gauss.high.high.identity} into \eqref{est_J11t4}, Bernstein's inequality and \eqref{sup3} yield
\begin{align*}
 &\sum_{\mu\ll\rho}
\sum_{\nu\sim\nu'\gtrsim\rho}
\rho\mu^{-1}
\sup_\tau
\left\|
\| \frac{\mu}{\nu} \cdot P_\mu L(P_\nu\lambda,P_{\nu'}\lambda)\|_{L_z^\infty}
\|\psi_{\sim\rho,1}\|_{L_z^2}
\|\psi_{\rho,1}^\tau\|_{L_z^2}
\right\|_{L_{t,y}^1}
\\
& \lesssim
\sum_{\mu\ll\rho}
\sum_{\nu\sim\nu'\gtrsim\rho}
\rho\mu^{d-3}\nu^{-1}
\|P_\nu\lambda\|_{L_{t,y}^4L_z^2} \cdot
\|P_{\nu'}\lambda\|_{L_{t,y}^4L_z^2}
 \cdot
\|\psi_{\sim\rho,1}\|_{L_{t,y}^4L_z^2}\cdot
\|\psi_{\rho,1}\|_{L_{t,y}^4L_z^2}
\\
& \lesssim
\varepsilon^{\frac32}
\rho^{-2s-2\sigma}
c_\rho^2(\sigma)
\sum_{\mu\ll\rho}
\left(\frac{\mu}{\rho}\right)^{d-3}\\
&\lesssim  \varepsilon^{\frac32}
\rho^{-2s-2\sigma}
c_\rho^2(\sigma),
\end{align*}
where \(d\geq5\),  and thus  the last sum is bounded.

Finally, the cubic terms we only need to consider \(L\left(
P_{\nu'}(\lambda  u_3),
\lambda_{\nu} \right)\) as follows, as the rest is the same.
\begin{align}
& \sum_{\mu\ll\rho}\sum_{ \nu\sim\nu'\gtrsim\rho }
\rho (\mu \nu)^{-1}
\Big\|
\|P_\mu
L\left(
P_{\nu'}(\lambda  u_3),
\lambda_{\nu}
\right)\|_{L_z^\infty}   
\|\psi_{\sim\rho,1}\|_{L_z^2}
\|\psi_{\rho,1} \|_{L_z^2}
\Big\|_{L_{t,y}^1} \label{cubic.high.high}\\
%&\lesssim	 \sum_{\mu\ll\rho}\sum_{ \nu\sim\nu'\gtrsim\rho }
%\rho\mu^{-1}\nu^{-1}\| P_\muL\left(P_{\nu'}(\lambda\cdot u_3),P_{\nu}\lambda\right)\|_{L^2_{t,y } L^\infty_z}  \cdot\|\psi_{\sim\rho,1}\|_{L_{t,y}^4L_z^2}\cdot\|\psi_{\rho,1}\|_{L_{t,y}^4L_z^2}\nonumber\\
&\lesssim	 \sum_{\mu\ll\rho}\sum_{ \nu\sim\nu'\gtrsim\rho }
\rho\mu^{d-\frac 52}\nu^{-1}\| P_\mu
L\left(
P_{\nu'}(\lambda\cdot u_3),
P_{\nu}\lambda
\right)\|_{L^2_{t } L^1_x}  \cdot
\|\psi_{\sim\rho,1}\|_{L_{t,y}^4L_z^2}\cdot
\|\psi_{\rho,1}\|_{L_{t,y}^4L_z^2} \nonumber\\
&\lesssim  	 \sum_{\mu\ll\rho}\sum_{ \nu\sim\nu'\gtrsim\rho }
\rho\mu^{d-\frac 52}\nu^{-1}\left\|P_{\nu^\prime } ( \lambda \cdot u_3) \right\|_{L^2_{t,x}}  \left\| P_{\nu'} \lambda   \right\|_{L^\infty_{t } L^2_x} \cdot
\|\psi_{\sim\rho,1}\|_{L_{t,y}^4L_z^2}\cdot
\|\psi_{\rho,1}\|_{L_{t,y}^4L_z^2}  \nonumber\\
&\lesssim   \sum_{\mu\ll\rho}\sum_{ \nu\sim\nu'\gtrsim\rho }
\rho\mu^{d-\frac 52}\nu^{-1} \varepsilon^{\frac34}(\nu')^{-s} c_{\nu'}(0)\cdot \varepsilon^{-\frac 18}\nu^{-s}c_{\nu}(0) \cdot  \varepsilon^{-\frac 14} \rho^{-2s+\frac12-2\sigma} c^{2}_{\rho} (\sigma) \nonumber\\
& \lesssim \varepsilon^{2}
\rho^{-2s-2\sigma}
c_\rho^2(\sigma). \nonumber
\end{align}

 

This completes the proof of this Lemma.

  
 \end{proof}

\end{document}
